\documentclass[../main.tex]{subfiles}

\begin{document}
% -----------------------------
\section{Batch}
% -----------------------------
A batch refers to processing multiple inputs together instead of executing a computation
separately for each input. In practice, it means running the same function or model multiple
times with different input values while treating them as a group.

Batch execution is useful because it allows repeated computations to be organized
and handled more efficiently. It simplifies workflows where the same model must be
evaluated on many inputs, such as testing, data processing, or machine learning inference.

Using batches is also convenient because it reduces repetitive code and provides
a clearer interface for running multiple evaluations, while often enabling better
performance in real systems that can process many inputs in parallel.

At the moment, we have a \path{run_module} that executes a single
function from the module using the interpreter.
It binds the provided runtime values to the function’s input arguments (entry
block arguments) and then interprets the operations in the block sequentially
until a \path{func.return} is reached, producing the final result.

\path{run_module_batch} extends this idea to batch execution. Instead of running
the function once, it repeatedly calls \path{run_module} with different argument sets,
allowing the same function to be evaluated multiple times with different inputs
and collecting all results into a list.

Start implementation by extending \texttt{Interpreter} class with \path{run_module_batch}
method:
\begin{lstlisting} [language=Python, caption={Batch execution implementation}]
    def run_module_batch(
        self,
        module,
        batch_args: list[list[int]],
        func_name: str = "my_func"
    ) -> list[int]:
        results = []

        for args in batch_args:
            # reset environment for each independent execution
            self.env = {}
            result = self.run_module(module, args=args, func_name=func_name)
            results.append(result)

        return results
\end{lstlisting}
Then we can modify our \path{hc_main.py} to use batch execution. Add a couple of input batches
and replace \path{hc_interpreter.run_module} with \path{hc_interpreter.run_module_batch} calls:
\begin{lstlisting} [language=Python, caption={Batch execution}]
    input_batches = [[65], [44], [40], [35], [30], [18]]
    ...
    result_before = hc_interpreter.run_module_batch(
        module, batch_args=input_batches, func_name="my_func"
    )
    ...
    result_after = lower_interpreter.run_module_batch(
        module, batch_args=input_batches, func_name="my_func"
    )
\end{lstlisting}

You should see score results spanning the range from the LO (0)
up to the HI parameter value (1000) defined in \path{build_model.py}.

\end{document}
